% from ega4-18-fr.tex
\oldpage[IV-4]{224}

\begin{corollary}[19.9.7]
\label{IV.19.9.7}
Sous les hypothèses générales de (19.9.6), la fonction $y\mapsto\operatorname{prof}_{Y_{f(y)},y}(\sh F_{f(y)})$ est semi-continue inférieurement dans $Y$. Si de plus $Y$ est localement constructible, cette fonction est localement constructible.
\end{corollary}

\begin{proposition}[19.9.8]
\label{IV.19.9.8}
Soient $f:X\to S$ un morphisme localement de présentation finie, $Z$ une partie fermée de $X$, $\sh F$ un $\sh O_X$-Module de présentation finie et $f$-plat, $\sh G$ un $\sh O_S$-Module quasi-cohérent. Supposons que, pour tout $x\in Z$, on ait $\operatorname{prof}((\sh F_{f(x)})_x)\geq1$ (resp. $\operatorname{prof}((\sh F_{f(x)})_x)\geq2$). Alors, si $j:X-Z\to X$ est l'injection canonique, et si on pose $\sh H=\sh F\otimes_{\sh O_X}f^*(\sh G)$, l'homomorphisme canonique $\rho_{\sh H}:\sh H\to j_*(j^*(\sh H))$ relatif à $j$ est injectif (resp. bijectif).
\end{proposition}

La question est évidemment locale sur $X$ et $S$, et il suffit de démontrer la proposition dans un voisinage d'un point $x\in Z$. Autrement dit, on peut se borner au cas où $X$ et $S$ sont affines. Grâce à (19.9.6), on peut supposer qu'il existe une section transversalement $\sh F$-régulière $g_1$ de $\sh O_X$ au-dessus de $X$ relativement à $S$ (resp. une suite transversalement $\sh F$-régulière $(g_1,g_2)$ de deux sections de $\sh O_X$ au-dessus de $X$ relativement à $S$) telle que si $Z'$ est l'ensemble des $x'\in X$ pour lesquels $g_1(x')=0$ (resp. $g_1(x')=g_2(x')=0$), on ait $Z\subset Z'$.

Cela étant, revenons à la démonstration de (19.9.8); en vertu de (19.9.6) on a encore $\operatorname{prof}((\sh F_{f(x)})_x)\geq1$ (resp. $\operatorname{prof}((\sh F_{f(x)})_x)\geq2$) pour tout $x\in Z'$. Si la proposition est prouvée en remplaçant le couple $(X,Z)$ par $(X,Z')$ et $(X-Z,Z'\cap(X-Z))$, il est immédiat qu'elle le sera aussi pour $(X,Z)$; on peut donc se borner à la prouver pour $X$ et $Z'$. Autrement dit, on peut se borner à démontrer (19.9.8) lorsque $S=\Spec(A)$, $X=\Spec(B)$, où $B$ est une $A$-algèbre de présentation finie, $Z=V(\mathfrak J)$, où $\mathfrak J$ est un idéal de type fini de l'anneau $B$, $\sh F=\widetilde M$ et $\sh G=\widetilde N$, où $N$ est un $A$-module et $M$ un $B$-module de présentation finie qui est un $A$-module plat. On peut en outre se ramener au cas où $N$ est un $A$-module de présentation finie. En effet, $N$ est limite inductive d'un système inductif filtrant de $A$-modules de présentation finie $N_\alpha$ (par double limite inductive) si l'on pose $\sh G_\alpha=\widetilde{N_\alpha}$, $\sh H$ est limite inductive des $\sh H_\alpha=\sh F\otimes_{\sh O_X}f^*(\sh G_\alpha)$; comme $j_*$ et $j^*$ commutent aux limites inductives et que $\varinjlim$ est un foncteur exact dans la catégorie des Modules quasi-cohérents, si la proposition est prouvée pour chacun des $\sh H_\alpha$, elle le sera pour $\sh H$.

Il suffit en outre de prouver que l'homomorphisme canonique
\[
  \Gamma(X,\sh H)\longrightarrow\Gamma(X-Z,\sh H)
\]
est injectif (resp. bijectif). Il existe alors un sous-anneau noethérien $A_0$ de $A$, une $A_0$-algèbre de type fini $B_0$, un idéal $\mathfrak J_0$ de $B_0$ et un $B_0$-module de type fini $M_0$ qui est

